\documentclass[11pt]{article}
\usepackage[margin=1in]{geometry}
\usepackage{amsmath,amssymb,booktabs,hyperref}
\title{Learned Epistemic Interrupts: Main-Model Control versus Semantic and Hidden-State Sidecars}
\author{Working Draft}\date{August 2026}
\begin{document}\maketitle
\begin{abstract}
Self-RAG already establishes learned on-demand retrieval and reflection in the
main generative LM \cite{selfrag}; learning when to retrieve is therefore not our
novelty claim. We instead test where epistemic-control competence should live:
inside the main LM, in a model-independent text sidecar observing ordinary
generation prefixes, or in a lightweight probe over a frozen LM's hidden state.
Controllers must choose among NONE, DB, LEXICAL, VECTOR, and WEB rather than
only retrieve/no-retrieve. We compare quality, unsupported claims, routing,
transfer, and cost without assuming that any architecture is superior.
\end{abstract}

\section{Problem Decomposition}
Separate:
\[
P(I_t\mid h_t)
\]
for whether retrieval is needed,
\[
P(S_t\mid I_t,h_t)
\]
for source selection, and
\[
P(Q_t\mid S_t,I_t,h_t)
\]
for normalized query formation. Detection, routing, and query formation must be evaluated independently.

\section{Candidate Interfaces}
Compare:
\begin{enumerate}
\item FLARE-like confidence control \cite{flare};
\item Self-RAG-style main-model learned retrieval/reflection control;
\item a text sidecar performing DETECT/ROUTE/NORMALIZE over ordinary prefixes;
\item a matched lightweight hidden-state sidecar with the main LM frozen;
\item oracle trigger, source, and query.
\end{enumerate}
If exact Self-RAG training is infeasible, the condition must be called a
Self-RAG-inspired approximation rather than Self-RAG proper. Main-model training
is confined to that declared comparison; text and hidden-state sidecars keep the
generative LM frozen.

\section{Training Data}
Use realistic generation prefixes, not only original user questions. Label:
\begin{itemize}
\item retrieve / do-not-retrieve;
\item information-need type;
\item preferred source;
\item canonical query/IR;
\item freshness/urgency;
\item evidence availability;
\item epistemic role.
\item miss-risk;
\item whether the need is computational rather than epistemic.
\end{itemize}

Hard negatives must include facts already present in context, quotations,
hypothetical or counterfactual statements, questions, user-given facts,
computational needs better served by calculators or logic, and fluent
non-factual prose.

\section{Epistemic Roles}
Distinguish imminent factual commitment, verification request, uncertainty,
quotation, hypothesis, counterfactual, user-provided fact, and computed or
derived fact. This prevents brittle surface matching from turning every
factual-looking phrase into a hard retrieval event.

\section{Routing}
Predict among
\[
\{\mathrm{NONE},\mathrm{DB},\mathrm{LEXICAL},\mathrm{VECTOR},\mathrm{WEB}\},
\]
with optional KG. Support abstention and, when justified, a small top-$k$ source
set. The DB, lexical, vector, and web actions must retain the distinct query and
execution semantics established in Paper 2.5.

\section{Sidecar Hypotheses and Ablations}
Model independence, upgradeability, routing specialization, small-model
leverage, and lower retraining cost are hypotheses rather than premises. Compare
confidence-only, semantic-only, and semantic-plus-confidence sidecars at matched
complexity and runtime cost. The hidden-state probe must also be complexity- and
cost-matched to the text controller; it earns no native-interface claim merely
by accessing activations.

For cross-model transfer, train the text sidecar on traces from one or several
model families and evaluate it on a different frozen family. Repeat across model
sizes. Successful transfer supports reusable runtime control; architecture- or
task-dependent winners are equally valid outcomes.

\section{Calibration and Cost}
False retrievals cost latency/money; missed retrievals cost unsupported claims. Evaluate thresholded policies under
\[
C=\lambda_{fp}C_{fp}+\lambda_{fn}C_{fn}+\lambda_{lat}C_{lat}+\lambda_{tok}C_{tok}.
\]
Do not optimize only F1.

\section{Experiments}
Evaluate in-domain detection, paraphrase robustness, unseen entities/relations,
domain and cross-model transfer, changing freshness windows, multi-factual
answers, and model-size effects. Measure the high-confidence/high-risk and
low-confidence/low-risk quadrants explicitly. Include DETECT/ROUTE/NORMALIZE,
confidence-feature, and action-space ablations.

\section{Baselines}
\begin{enumerate}
\item no retrieval;
\item upfront RAG;
\item explicit search/tool calls;
\item Paper 2.5 heuristic reflex routing;
\item FLARE-like confidence controller;
\item Self-RAG-style main-model learned controller;
\item text sidecar DETECT/ROUTE/NORMALIZE;
\item frozen-LM hidden-state sidecar;
\item oracle trigger/source/query.
\end{enumerate}

\section{Metrics}
Report trigger AUROC/AUPRC and P/R, high-confidence/high-risk recall,
low-confidence/low-risk false retrieval, calibration overall and by source,
selective-risk/coverage curves, routing top-$1$/top-$k$ and confusion matrices,
normalization correctness, unsupported claims, final accuracy, count/fanout,
latency/cost, model-size interactions, OOD and cross-model degradation, and the
cost versus unsupported-claim Pareto frontier.

\section{Falsification}
The sidecar architecture is weakened if main-model Self-RAG-style control
dominates quality and cost, sidecars do not transfer, heterogeneous routing adds
no value, confidence explains semantic gains, hidden-state dependence is too
model-specific, or the detector memorizes surface phrases. A main-model,
sidecar, hybrid, or task/model-size-dependent winner is a valid result.

\section{Scope}
No joint base-model training except the declared main-model comparison; no
transactional or backtracking architecture as the main contribution; and no
PRA or native KV dependency. Relevance and support labels are prior art, not a novelty
claim. Paper 4 concerns persistent cross-source provenance, dependency
propagation, transactional scopes, retraction, and backtracking.

\begin{thebibliography}{9}
\bibitem{flare}
Z.~Jiang et al.
Active retrieval augmented generation.
\emph{EMNLP}, 2023.
\url{https://arxiv.org/abs/2305.06983}.

\bibitem{selfrag}
A.~Asai et al.
Self-RAG: Learning to retrieve, generate, and critique through self-reflection.
\emph{ICLR}, 2024.
\url{https://arxiv.org/abs/2310.11511}.
\end{thebibliography}

\end{document}
